% §11.7 积分在物理中的应用

\section{积分在物理中的应用}

在场论这一节, 我们知道曲线积分, 曲面积分以及 Green 公式, Stokes 公式, Gauss 公式都对物理学有着特别重要的意义. 现在将给出各种各样的具体实例.

\subsection*{1. 物体的质量}

设 $C\subset\mathbb{R}^{3}$ 为可求长的曲线, 其密度函数 $\rho:C\to\mathbb{R}$ 是连续的. 如果 $C$ 分为 $m$ 个小段 $C_{1},C_{2},\cdots,C_{m}$, 每一小段 $C_{i}$ 上取一点 $\boldsymbol{\xi}^{i}$, 其质量近似为 $\rho(\boldsymbol{\xi}^{i})\Delta s_{i}$, 其中 $\Delta s_{i}$ 为 $C_{i}$ 的弧长. 于是,
\[
\sum_{i=1}^{m}\rho(\boldsymbol{\xi}^{i})\Delta s_{i}
\]
可以作为 $C$ 的质量的近似值. 自然,
\[
\int_{C}\rho\,\mathrm{d}s=\lim_{\max\Delta s_{i}\to0}\sum_{i=1}^{m}\rho(\boldsymbol{\xi}^{i})\Delta s_{i}
\]
就是曲线 $C$ 的质量的精确值. 这是第一型曲线积分.

如果 $\Sigma\subset\mathbb{R}^{3}$ 为有面积的二维曲面, 其密度函数 $\rho:\Sigma\to\mathbb{R}$ 是连续的. 类似上述极限过程 (以后不再赘述), 曲面 $\Sigma$ 的质量为
\[
\iint_{\Sigma}\rho\,\mathrm{d}\sigma.
\]
这是第一型曲面积分.

如果 $\Omega\subset\mathbb{R}^{3}$ 为有体积的立体, 其密度函数 $\rho:\Omega\to\mathbb{R}$ 是连续的, 则 $\Omega$ 的质量为
\[
\iiint_{\Omega}\rho\,\mathrm{d}v.
\]
这是三重积分, 是第一型曲面积分的特殊情形.

注意, 物体是什么形状 (曲线, 曲面, 立体), 就决定了采用什么类型的积分 (第一型曲线积分, 第一型曲面积分, 三重积分).

\subsection*{2. 物体的质心}

设 $C\subset\mathbb{R}^{3}$ 为可求长曲线, 其密度函数 $\rho:C\to\mathbb{R}$ 是连续的, 则 $C$ 的质心为
\[
(x_{0},y_{0},z_{0})=\frac{1}{\int_{C}\rho\,\mathrm{d}s}\Bigl(\int_{C}x\rho\,\mathrm{d}s,\ \int_{C}y\rho\,\mathrm{d}s,\ \int_{C}z\rho\,\mathrm{d}s\Bigr);
\]
当 $C$ 是均匀时, 令 $\rho=1$.

设 $\Sigma\subset\mathbb{R}^{3}$ 为有面积的二维曲面时, 其密度函数 $\rho:\Sigma\to\mathbb{R}$ 是连续的, 则 $\Sigma$ 的质心为
\[
(x_{0},y_{0},z_{0})=\frac{1}{\iint_{\Sigma}\rho\,\mathrm{d}\sigma}\Bigl(\iint_{\Sigma}x\rho\,\mathrm{d}\sigma,\ \iint_{\Sigma}y\rho\,\mathrm{d}\sigma,\ \iint_{\Sigma}z\rho\,\mathrm{d}\sigma\Bigr);
\]
当 $\Sigma$ 是均匀时, 令 $\rho=1$.

设 $\Omega\subset\mathbb{R}^{3}$ 为有体积的立体, 其密度函数 $\rho:\Omega\to\mathbb{R}$ 是连续的, 则 $\Omega$ 的质心为
\[
(x_{0},y_{0},z_{0})=\frac{1}{\iiint_{\Omega}\rho\,\mathrm{d}v}\Bigl(\iiint_{\Omega}x\rho\,\mathrm{d}v,\ \iiint_{\Omega}y\rho\,\mathrm{d}v,\ \iiint_{\Omega}z\rho\,\mathrm{d}v\Bigr);
\]
当 $\Omega$ 是均匀时, 令 $\rho=1$.

在平面上, 密度为 $\rho$ (连续)的曲线 $C$ 的质心为
\[
(x_{0},y_{0})=\frac{1}{\int_{C}\rho\,\mathrm{d}s}\Bigl(\int_{C}x\rho\,\mathrm{d}s,\ \int_{C}y\rho\,\mathrm{d}s\Bigr);
\]
密度为 $\rho$ (连续)的有面积集 $\Sigma$ 的质心为
\[
(x_{0},y_{0})=\frac{1}{\iint_{\Sigma}\rho\,\mathrm{d}x\mathrm{d}y}\Bigl(\iint_{\Sigma}x\rho\,\mathrm{d}x\mathrm{d}y,\ \iint_{\Sigma}y\rho\,\mathrm{d}x\mathrm{d}y\Bigr).
\]

\subsection*{3. 物体的转动惯量}

在平面上, 密度为 $\rho$ (连续)的曲线 $C$ 对于 $Ox$ 轴与 $Oy$ 轴的转动惯量分别为
\[
I_{x}=\int_{C}y^{2}\rho\,\mathrm{d}s,\quad I_{y}=\int_{C}x^{2}\rho\,\mathrm{d}s.
\]
密度为 $\rho$ (连续)的有面积集 $\Sigma$ 对于 $Ox$ 轴与 $Oy$ 轴的转动惯量分别为
\[
I_{x}=\iint_{\Sigma}y^{2}\rho\,\mathrm{d}x\mathrm{d}y,\quad I_{y}=\iint_{\Sigma}x^{2}\rho\,\mathrm{d}x\mathrm{d}y.
\]

在三维空间中, 设 $\Omega\subset\mathbb{R}^{3}$ 为有体积的立体, 其密度函数 $\rho:\Omega\to\mathbb{R}$ 是连续的. 三重积分
\[
\begin{aligned}
I_{xy}&=\iiint_{\Omega}z^{2}\rho\,\mathrm{d}x\mathrm{d}y\mathrm{d}z,\quad I_{yz}=\iiint_{\Omega}x^{2}\rho\,\mathrm{d}x\mathrm{d}y\mathrm{d}z,\\
I_{zx}&=\iiint_{\Omega}y^{2}\rho\,\mathrm{d}x\mathrm{d}y\mathrm{d}z
\end{aligned}
\]
分别称为物体 $\Omega$ 对于坐标平面 $xOy$, $yOz$, $zOx$ 的转动惯量.

积分
\[
I_{l}=\iiint_{\Omega}r^{2}\rho\,\mathrm{d}x\mathrm{d}y\mathrm{d}z
\]
(其中 $r$ 为物体 $\Omega$ 的动点 $(x,y,z)$ 与轴 $l$ 的距离)称为物体 $\Omega$ 对于轴 $l$ 的转动惯量. 特别是, 对于坐标轴 $Ox$, $Oy$, $Oz$ 分别有
\[
I_{x}=I_{xy}+I_{xz},\quad I_{y}=I_{yx}+I_{yz},\quad I_{z}=I_{zx}+I_{zy}.
\]
积分
\[
I_{O}=\iiint_{\Omega}(x^{2}+y^{2}+z^{2})\rho\,\mathrm{d}x\mathrm{d}y\mathrm{d}z=I_{xy}+I_{yz}+I_{zx}
\]
称为物体 $\Omega$ 对坐标原点的转动惯量.

关于 $\mathbb{R}^{3}$ 中密度为 $\rho$ (连续)的可求长曲线 $C$ 与密度为 $\rho$ (连续)的有面积的曲面, 可类似定义各种转动惯量. 前者是第一型曲线积分, 后者是第一型曲面积分.

当 $\rho=1$ 时, 称为几何转动惯量.

\subsection*{4. 引力场的位}

积分
\[
u(x,y,z)=\iiint_{\Omega}\rho(\xi,\eta,\zeta)\frac{\mathrm{d}\xi\mathrm{d}\eta\mathrm{d}\zeta}{r}
\]
称为物体 $\Omega$ 在点 $(x,y,z)$ 的 Newton 位, 其中 $\Omega\subseteq\mathbb{R}^{3}$ 为有体积的立体, $\rho:\Omega\to\mathbb{R}$ 为 $\Omega$ 的密度函数, $\rho$ 连续, $r=\sqrt{(\xi-x)^{2}+(\eta-y)^{2}+(\zeta-z)^{2}}$.

质量为 $m$ 的质点吸引物体的力在坐标轴 $Ox$, $Oy$, $Oz$ 轴上的投影 $X$, $Y$, $Z$ 分别为
\[
X=km\frac{\partial u}{\partial x}=km\iiint_{\Omega}\rho\frac{\xi-x}{r^{3}}\,\mathrm{d}\xi\mathrm{d}\eta\mathrm{d}\zeta,
\]
\[
Y=km\frac{\partial u}{\partial y}=km\iiint_{\Omega}\rho\frac{\eta-y}{r^{3}}\,\mathrm{d}\xi\mathrm{d}\eta\mathrm{d}\zeta,
\]
\[
Z=km\frac{\partial u}{\partial z}=km\iiint_{\Omega}\rho\frac{\zeta-z}{r^{3}}\,\mathrm{d}\xi\mathrm{d}\eta\mathrm{d}\zeta,
\]
其中 $k$ 为引力定律常数.

\subsection*{5. 变力做功、环量与涡量}

称第二型曲线积分
\[
\begin{aligned}
\int_{\overrightarrow{C}}\boldsymbol{F}\cdot\mathrm{d}\boldsymbol{x}&=\int_{\overrightarrow{C}}P(x,y,z)\mathrm{d}x+Q(x,y,z)\mathrm{d}y+R(x,y,z)\mathrm{d}z\\
&=\int_{\alpha}^{\beta}\bigl[P(x(t),y(t),z(t))x'(t)+Q(x(t),y(t),z(t))y'(t)\\
&\quad+R(x(t),y(t),z(t))z'(t)\bigr]\mathrm{d}t
\end{aligned}
\]
为变力 $\boldsymbol{F}=(P,Q,R)$ 沿定向曲线 $\overrightarrow{C}$ 所做的功.

当 $\overrightarrow{C}$ 为一条 $C^{1}$ 的定向封闭曲线时, 称 $\int_{\overrightarrow{C}}\boldsymbol{F}\cdot\mathrm{d}\boldsymbol{x}$ 为向量场 $\boldsymbol{F}$ 沿定向曲线 $\overrightarrow{C}$ 的环量. 若 $\overrightarrow{C}$ 为有向曲面 $\overrightarrow{M}$ 的边界时, $\overrightarrow{C}$ 的方向为 $\overrightarrow{M}$ 的诱导方向, 则有
\[
\oint_{\overrightarrow{\partial M}}\boldsymbol{F}\cdot\boldsymbol{\tau}\,\mathrm{d}s=\oint_{\overrightarrow{\partial M}}\boldsymbol{F}\cdot\mathrm{d}\boldsymbol{x}\reason{Stokes 公式}\iint_{\overrightarrow{M}}\mathrm{rot}\,\boldsymbol{F}\cdot\boldsymbol{n}_{0}\,\mathrm{d}\sigma,
\]
其中 $\boldsymbol{\tau}$ 为与 $\overrightarrow{C}$ 方向一致的连续单位切向量场, $\boldsymbol{n}_{0}$ 为与定向曲面 $\overrightarrow{M}$ 方向一致的连续单位法向量场, $\boldsymbol{F}$ 为 $C^{r}\ (r\geqslant1)$ 向量场.

设 $\boldsymbol{F}=(P,Q,R)$ 为开区域 $U\subset\mathbb{R}^{3}$ 中的 $C^{r}\ (r\geqslant1)$ 流速场, $\boldsymbol{n}_{0}(\boldsymbol{p})$ 为 $\boldsymbol{p}\in U$ 处的单位法向量, 作垂直于 $\boldsymbol{n}_{0}(\boldsymbol{p})$ 的圆盘 $M$, 则
\[
\oint_{\overrightarrow{\partial M}}\boldsymbol{F}\cdot\boldsymbol{\tau}\,\mathrm{d}s=\oint_{\overrightarrow{\partial M}}\boldsymbol{F}\cdot\mathrm{d}\boldsymbol{x}
\]
反映了流体环绕着圆周 $\partial M$ 的旋转程度, 而
\[
\frac{1}{\sigma(M)}\oint_{\overrightarrow{\partial M}}\boldsymbol{F}\cdot\mathrm{d}\boldsymbol{x}
\]
为平均旋转强度. 保持方向 $\boldsymbol{n}_{0}$ 不变, 令 $M\to\boldsymbol{p}$ ($M$ 缩成点 $\boldsymbol{p}$) 得到
\[
\begin{aligned}
\lim_{M\to\boldsymbol{p}}\frac{1}{\sigma(M)}\oint_{\overrightarrow{\partial M}}\boldsymbol{F}\cdot\mathrm{d}\boldsymbol{x}&\reason{Stokes 公式}\lim_{M\to\boldsymbol{p}}\frac{1}{\sigma(M)}\iint_{\overrightarrow{M}}\mathrm{rot}\,\boldsymbol{F}\cdot\boldsymbol{n}_{0}\,\mathrm{d}\sigma\\
&\reason{积分中值定理}\lim_{M\to\boldsymbol{p}}\frac{1}{\sigma(M)}\mathrm{rot}\,\boldsymbol{F}(\boldsymbol{\xi})\cdot\boldsymbol{n}_{0}\iint_{\overrightarrow{M}}\mathrm{d}\sigma\\
&=\mathrm{rot}\,\boldsymbol{F}(\boldsymbol{p})\cdot\boldsymbol{n}_{0},
\end{aligned}
\]
它是 $\mathrm{rot}\,\boldsymbol{F}(\boldsymbol{p})$ 在方向 $\boldsymbol{n}_{0}$ 上的投影, 是流速场 $\boldsymbol{F}$ 在 $\boldsymbol{p}$ 点关于方向 $\boldsymbol{n}_{0}$ 的旋转强度, 也称为流体在 $\boldsymbol{p}$ 点处绕 $\boldsymbol{n}_{0}$ 方向的涡量.

当 $\mathrm{rot}\,\boldsymbol{F}$ 与 $\boldsymbol{n}_{0}$ 同向时, 绕 $\boldsymbol{n}_{0}$ 的涡量最大, 其值为 $\|\mathrm{rot}\,\boldsymbol{F}(\boldsymbol{p})\|$.

\subsection*{6. 通量}

设 $M\subseteq\mathbb{R}^{3}$ 为一个区域, $\overrightarrow{M}$ 取右手系, $\overrightarrow{\partial M}$ 为取外法向的定向封闭曲面, $\boldsymbol{n}_{0}$ 为与 $\overrightarrow{\partial M}$ 的方向一致的连续的单位法向量. $\boldsymbol{F}=(P,Q,R)$ 为 $C^{1}$ 流速场, 则有 Gauss 公式
\[
\iint_{\overrightarrow{\partial M}}\boldsymbol{F}\cdot\boldsymbol{n}_{0}\,\mathrm{d}\sigma=\iiint_{\overrightarrow{M}}\mathrm{div}\,\boldsymbol{F}\,\mathrm{d}v.
\]
上式左边的 $\boldsymbol{F}\cdot\boldsymbol{n}_{0}\mathrm{d}\sigma$ 为单位时间内流过面积为 $\mathrm{d}\sigma$ 的小曲面的流量, 而 $\oint_{\overrightarrow{\partial M}}\boldsymbol{F}\cdot\boldsymbol{n}_{0}\,\mathrm{d}\sigma$ 为单位时间内流出定向封闭曲面 $\overrightarrow{\partial M}$ 的总量, 称为通量. 这是一个代数和. 称
\[
\frac{1}{v(M)}\oint_{\overrightarrow{\partial M}}\boldsymbol{F}\cdot\boldsymbol{n}_{0}\,\mathrm{d}\sigma
\]
为单位体积内流源的「平均散射强度」.

如果
\[
\oint_{\overrightarrow{\partial M}}\boldsymbol{F}\cdot\boldsymbol{n}_{0}\,\mathrm{d}\sigma>0,
\]
$M$ 内有「源」(图 \ref{fig:11.7.1}), 表明从 $\overrightarrow{\partial M}$ 内流出; 如果
\[
\oint_{\overrightarrow{\partial M}}\boldsymbol{F}\cdot\boldsymbol{n}_{0}\,\mathrm{d}\sigma<0,
\]
$M$ 内有「汇」或「负源」(图 \ref{fig:11.7.2}), 表明向 $\overrightarrow{\partial M}$ 内流进.

\begin{figure}[htbp]
  \centering
  \begin{minipage}[b]{0.46\linewidth}
    \centering
    \begin{tikzpicture}[scale=1.05, >=stealth, line cap=round, line join=round]
      \draw[thick, fill=ThemeColor!6] (0,0) circle (1.35);
      \foreach \a in {0, 45, 90, 135, 180, 225, 270, 315} {
        \draw[->, very thick, ThemeColor] ({0.4*cos(\a)}, {0.4*sin(\a)}) -- ({1.85*cos(\a)}, {1.85*sin(\a)});
      }
      \fill (0,0) circle (1.6pt) node[above right] {源};
      \node[below] at (0,-1.55) {$\vec{\partial M}$ 内流出};
    \end{tikzpicture}
    \caption{源}\label{fig:11.7.1}
  \end{minipage}\hfill
  \begin{minipage}[b]{0.46\linewidth}
    \centering
    \begin{tikzpicture}[scale=1.05, >=stealth, line cap=round, line join=round]
      \draw[thick, fill=ThemeColor!6] (0,0) circle (1.35);
      \foreach \a in {0, 45, 90, 135, 180, 225, 270, 315} {
        \draw[->, very thick, ThemeColor] ({1.85*cos(\a)}, {1.85*sin(\a)}) -- ({0.4*cos(\a)}, {0.4*sin(\a)});
      }
      \fill (0,0) circle (1.6pt) node[above right] {汇};
      \node[below] at (0,-1.55) {向 $\vec{\partial M}$ 内流进};
    \end{tikzpicture}
    \caption{汇}\label{fig:11.7.2}
  \end{minipage}
\end{figure}

如果
\[
\oint_{\overrightarrow{\partial M}}\boldsymbol{F}\cdot\boldsymbol{n}_{0}\,\mathrm{d}\sigma=0,
\]
汇与源相抵消. 如果 $\mathrm{div}\,\boldsymbol{F}=0$, 则称场 $\boldsymbol{F}$ 为无源场.

在 $M$ 内任取一个内点 $\boldsymbol{p}_{0}$, 以 $\boldsymbol{p}_{0}$ 为球心, 作一个半径为 $\varepsilon$ 的小球 $B(\boldsymbol{p}_{0};\varepsilon)$, 让 $\varepsilon$ 如此之小, 使得 $B(\boldsymbol{p}_{0};\varepsilon)\subset M$. $B(\boldsymbol{p}_{0};\varepsilon)$ 的边界 $\partial B(\boldsymbol{p}_{0};\varepsilon)$ 为小球的球面, 取 $\boldsymbol{n}_{0}$ 为该小球面的连续的外单位法向量场. 我们有
\[
\begin{aligned}
\iint_{\partial B(\boldsymbol{p}_{0};\varepsilon)}\boldsymbol{F}\cdot\boldsymbol{n}_{0}\,\mathrm{d}\sigma&\reason{Gauss 公式}\iiint_{B(\boldsymbol{p}_{0};\varepsilon)}\mathrm{div}\,\boldsymbol{F}\,\mathrm{d}v\\
&\reasonb{积分中值定理}{\boldsymbol{\xi}\in\overline{B(\boldsymbol{p}_{0};\varepsilon)}}\mathrm{div}\,\boldsymbol{F}(\boldsymbol{\xi})v(B(\boldsymbol{p}_{0};\varepsilon)),
\end{aligned}
\]
其中 $B(\boldsymbol{p}_{0};\varepsilon)$ 为右手系. 由此得到
\[
\mathrm{div}\,\boldsymbol{F}(\boldsymbol{\xi})=\frac{1}{v(B(\boldsymbol{p}_{0};\varepsilon))}\iint_{\partial B(\boldsymbol{p}_{0};\varepsilon)}\boldsymbol{F}\cdot\boldsymbol{n}_{0}\,\mathrm{d}\sigma,
\]
这是单位体积上流速场送出的流体量. 如果令 $\varepsilon\to0^{+}$, 即让小球无限地收缩到 $\boldsymbol{p}_{0}$ 点时, 立即得到
\[
\mathrm{div}\,\boldsymbol{F}(\boldsymbol{p}_{0})=\lim_{\varepsilon\to0^{+}}\mathrm{div}\,\boldsymbol{F}(\boldsymbol{\xi})=\lim_{\varepsilon\to0^{+}}\frac{1}{v(B(\boldsymbol{p}_{0};\varepsilon))}\iint_{\partial B(\boldsymbol{p}_{0};\varepsilon)}\boldsymbol{F}\cdot\boldsymbol{n}_{0}\,\mathrm{d}\sigma.
\]
由此可见, 散度刻画在一点处流速场产生流体的能力. 它是流源在点 $\boldsymbol{p}_{0}$ 的散射强度, 称为 $\boldsymbol{F}$ 在 $\boldsymbol{p}_{0}$ 点散度. 如果 $\mathrm{div}\,\boldsymbol{F}=0$, 称 $\boldsymbol{F}$ 为无源场, 此时, 由 Gauss 公式知, 流过任何定向封闭曲面的通量恒为 $0$; 反之, 由上面公式推得 $\mathrm{div}\,\boldsymbol{F}=0$.

类似可讨论磁场, 磁通量.

\begin{example}\label{ex:11.7.1}
静电场的 Gauss 定理: 在静电场中, 通过任一封闭曲面的电通量等于此曲面所包含的电荷总量的 $4\pi$ 倍.
\end{example}

\begin{proof}
设 $\Sigma$ 为一封闭曲面, 取单位外法向 $\boldsymbol{n}_{0}$ 的方向为其定向. $\Sigma$ 所包围的定向区域为 $\overrightarrow{M}$, 它的方向为右手系, 则 $\overrightarrow{\Sigma}=\overrightarrow{\partial M}$. 分三步来证.

(1) 在坐标原点上, 放置电量为 $q$ 的点电荷. 这个点电荷产生的电场强度 (简称场强)是向量
\[
\boldsymbol{E}(\boldsymbol{p})=\frac{q}{p^{3}}\boldsymbol{p},\quad\boldsymbol{p}\neq\boldsymbol{0},
\]
其中 $p=\|\boldsymbol{p}\|$. 由例 \ref{ex:11.6.1}(2) 知, 当 $\boldsymbol{p}\neq\boldsymbol{0}$ 时,
\[
\mathrm{div}\,\boldsymbol{E}=0.
\]
如果原点 $O$ 在 $\Sigma$ 所围成的体 $M$ 的外部, 则 $\boldsymbol{E}$ 通过定向封闭曲面 $\Sigma$ 的通量为零, 即
\[
\iint_{\Sigma}\boldsymbol{E}\cdot\boldsymbol{n}_{0}\,\mathrm{d}\sigma\reason{Gauss}\iiint_{M}\mathrm{div}\,\boldsymbol{E}\,\mathrm{d}v=\iiint_{M}0\,\mathrm{d}v=0.
\]
再设 $O$ 被曲面 $\Sigma$ 包围在内部. 这时, 以 $O$ 为中心, 充分小的 $\varepsilon>0$ 为半径作一小球面 $\Sigma_{\varepsilon}$, 使得这个小球面完全落在曲面 $\Sigma$ 的包围之内, 由 $\Sigma$ 与 $\Sigma_{\varepsilon}$ 所围成的区域 $M_{\varepsilon}$ 取右手系得到定向区域 $\overrightarrow{M}_{\varepsilon}$. 记 $\overrightarrow{\Sigma}_{\varepsilon}$ 为单位外法向 $\boldsymbol{n}_{0}$ 的定向球面 (图 \ref{fig:11.7.3}), 则 $\overrightarrow{\Sigma}+\overrightarrow{\Sigma}_{\varepsilon}^{-}=\overrightarrow{\partial M}_{\varepsilon}$ ($\overrightarrow{M}_{\varepsilon}$ 的诱导定向). 于是,
\[
\iint_{\Sigma}\boldsymbol{E}\cdot\boldsymbol{n}_{0}\,\mathrm{d}\sigma+\iint_{\Sigma_{\varepsilon}}\boldsymbol{E}\cdot(-\boldsymbol{n}_{0})\,\mathrm{d}\sigma=\iiint_{M_{\varepsilon}}\mathrm{div}\,\boldsymbol{E}\,\mathrm{d}v=\iiint_{M_{\varepsilon}}0\,\mathrm{d}v=0,
\]

\begin{figure}[htbp]
  \centering
  \begin{tikzpicture}[scale=1.1, >=stealth, line cap=round, line join=round]
    % 外曲面 Sigma
    \draw[thick, fill=ThemeColor!6, even odd rule]
      plot[smooth cycle, tension=0.88] coordinates {(0,0) (1.7,-0.3) (3.0,0.6) (3.2,2.0) (2.2,3.15) (0.9,3.25) (-0.3,2.3) (-0.5,1.0)}
      (1.3,1.55) circle (0.65);
    \node at (3.3,2.15) {$\Sigma$};
    
    % 内小球面 Sigma_epsilon
    \draw[thick, ThemeColor] (1.3,1.55) circle (0.65);
    \fill (1.3,1.55) circle (1.5pt) node[left] {$O$};
    \node[below right, ThemeColor] at (1.75,1.1) {$\Sigma_{\varepsilon}$};
    
    % 外法向量
    \draw[->, very thick, ThemeColor] (2.9,0.7) -- (3.5,0.7) node[right] {$\boldsymbol{n}$};
    \draw[->, very thick, ThemeColor] (1.75,1.95) -- (2.05,2.25) node[above right] {$\boldsymbol{n}_0$};
    \node at (1.8,2.7) {$M_\varepsilon$};
  \end{tikzpicture}
  \caption{静电场的 Gauss 定理}\label{fig:11.7.3}
\end{figure}

\[
\begin{aligned}
\iint_{\Sigma}\boldsymbol{E}\cdot\boldsymbol{n}_{0}\,\mathrm{d}\sigma&=\iint_{\Sigma_{\varepsilon}}\boldsymbol{E}\cdot\boldsymbol{n}_{0}\,\mathrm{d}\sigma=\iint_{\Sigma_{\varepsilon}}\left(\frac{q}{p^{3}}\boldsymbol{p}\right)\cdot\frac{\boldsymbol{p}}{p}\,\mathrm{d}\sigma\\
&=\iint_{\Sigma_{\varepsilon}}\frac{q}{\varepsilon^{2}}\,\mathrm{d}\sigma=\frac{q}{\varepsilon^{2}}4\pi\varepsilon^{2}=4\pi q.
\end{aligned}
\]

(2) 设想有限个点电荷, 分别带电量 $q_{1},q_{2},\cdots,q_{k}$. 又设点电荷 $q_{i}$ 产生的场强为 $\boldsymbol{E}_{i}$ ($i=1,2,\cdots,k$). 因此, 它们产生的总场强
\[
\boldsymbol{E}=\boldsymbol{E}_{1}+\boldsymbol{E}_{2}+\cdots+\boldsymbol{E}_{k}.
\]
如果点电荷 $q_{1},q_{2},\cdots,q_{t}$ 被包含在 $\Sigma$ 的内部, 其他的点电荷 $q_{t+1},\cdots,q_{k}$ 在 $\Sigma$ 的外部, 由 (1) 可知
\[
\iint_{\Sigma}\boldsymbol{E}_{i}\cdot\boldsymbol{n}_{0}\,\mathrm{d}\sigma=\begin{cases}
4\pi q_{i}, & i=1,2,\cdots,t,\\
0, & i=t+1,\cdots,k.
\end{cases}
\]
因此,
\[
\iint_{\Sigma}\boldsymbol{E}\cdot\boldsymbol{n}_{0}\,\mathrm{d}\sigma=\sum_{i=1}^{k}\iint_{\Sigma}\boldsymbol{E}_{i}\cdot\boldsymbol{n}_{0}\,\mathrm{d}\sigma=4\pi(q_{1}+q_{2}+\cdots+q_{t}).
\]

(3) 设区域 $U\subset\mathbb{R}^{3}$ 内的场强 $\boldsymbol{E}$ 是由连续分布的电荷所产生的, 电荷密度为 $\rho(\boldsymbol{p})$, 它是点 $\boldsymbol{p}$ 的连续函数. 设 $M\subset U$ 为区域, 其内包含的总电荷为三重积分 $\iiint_{M}\rho\,\mathrm{d}v$. 于是
\[
\iiint_{M}\mathrm{div}\,\boldsymbol{E}\,\mathrm{d}v=\iint_{\partial M}\boldsymbol{E}\cdot\boldsymbol{n}_{0}\,\mathrm{d}\sigma=4\pi\iiint_{M}\rho\,\mathrm{d}v,
\]
即
\[
\iiint_{M}(\mathrm{div}\,\boldsymbol{E}-4\pi\rho)\,\mathrm{d}v=0.
\]
因为 $M$ 可以是 $U$ 中任何立体, 所以由上式及反证法可知
\[
\mathrm{div}\,\boldsymbol{E}-4\pi\rho=0
\]
在 $U$ 上处处成立. 于是,
\[
\mathrm{div}\,\boldsymbol{E}=4\pi\rho.
\]
这是静电场的基本公式之一.
\end{proof}

\begin{example}\label{ex:11.7.2}
设 $\boldsymbol{F}$ 为流速场, 数量场 $\rho$ 为流体在各点处的密度. 场 $\boldsymbol{F}$ 和 $\rho$ 既依赖于空间点的位置, 也依赖于时间 $t$. 又设 $M$ 为空间中任意固定的区域. 则流体的质量关于时间 $t$ 的变化率应该为
\[
\frac{\mathrm{d}}{\mathrm{d}t}\iiint_{M}\rho(\boldsymbol{p},t)\,\mathrm{d}v=\iiint_{M}\frac{\partial\rho}{\partial t}(\boldsymbol{p},t)\,\mathrm{d}v.
\]
由于流体是不可压缩的, 上述变化率必须等于流体进入 $M$ 的速率, 即
\[
\frac{\mathrm{d}}{\mathrm{d}t}\iiint_{M}\rho(\boldsymbol{p},t)\,\mathrm{d}v=-\oint_{\partial M}\rho\boldsymbol{F}\cdot\boldsymbol{n}_{0}\,\mathrm{d}\sigma,
\]
于是,
\[
\begin{aligned}
\iiint_{M}\frac{\partial\rho}{\partial t}(\boldsymbol{p},t)\,\mathrm{d}v&=\frac{\mathrm{d}}{\mathrm{d}t}\iiint_{M}\rho(\boldsymbol{p},t)\,\mathrm{d}v=-\iint_{\partial M}\rho\boldsymbol{F}\cdot\boldsymbol{n}_{0}\,\mathrm{d}\sigma\\
&\reason{Gauss 公式}-\iiint_{M}\mathrm{div}(\rho\boldsymbol{F})\,\mathrm{d}v,
\end{aligned}
\]
\[
\iiint_{M}\bigl[\frac{\partial\rho}{\partial t}+\mathrm{div}(\rho\boldsymbol{F})\bigr]\mathrm{d}v=0.
\]
由于被积函数的连续性, $M$ 是任意区域及反证法可知
\[
\frac{\partial\rho}{\partial t}+\mathrm{div}(\rho\boldsymbol{F})=0
\]
在 $\mathbb{R}^{3}$ 中处处成立, 这个方程称为流体的连续性方程.
\end{example}

\begin{exblock}{练习题 11.7}
\item 定向曲线 $\overrightarrow{C}$: $x^{2}+y^{2}+z=a^{2}$, $x^{2}+y^{2}=ax$, $z\geqslant0\ (a>0)$, 从点 $\left(\frac{a}{2},0,0\right)$ 望去, 沿逆时针方向行进. 计算力场 $\boldsymbol{F}=(y^{2},z^{2},x^{2})$ 沿 $\overrightarrow{C}$ 所做的功.

\item 质量为 $m$ 的质点在力场 $\boldsymbol{F}$ 的作用下沿曲线 $\overrightarrow{C}$ 运动, $\overrightarrow{C}$ 的起点为 $\boldsymbol{a}$, 终点为 $\boldsymbol{b}$. 用 $\boldsymbol{v}$ 记质点的速度向量. 证明: 力场 $\boldsymbol{F}$ 所做的功
\[
\int_{\overrightarrow{C}}\boldsymbol{F}\cdot\mathrm{d}\boldsymbol{x}=\frac{1}{2}m\boldsymbol{v}^{2}(\boldsymbol{p})\Big|_{\boldsymbol{a}}^{\boldsymbol{b}}.
\]

\item 设有流速场 $\boldsymbol{F}=(yz,zx,xy)$, 曲面 $\Sigma$ 是圆柱体 $x^{2}+y^{2}\leqslant a^{2}$, $0\leqslant z\leqslant h$ 的表面, 求流速场 $\boldsymbol{F}$ 流出 $\Sigma$ 的流量.

\item 给定流速场 $\boldsymbol{F}=(y,z,x)$, 封闭曲面 $\Sigma:x^{2}+y^{2}=R^{2}\ (R>0)$, $z=0$, $z=h$. 计算 $\boldsymbol{F}$ 流向曲面 $\Sigma$ 之外的流量.

\item 设有向曲面 $\Sigma$ 为
\[
x+y+z=1,\quad x,y,z\geqslant0,
\]
其法向与向量 $(1,1,1)$ 同向, 求力场 $\boldsymbol{F}=(y^{2},z^{2},x^{2})$ 绕 $\Sigma$ 的诱导定向的边界 $\partial\Sigma$ 一周所做的功.

\item 弹性力的方向指向坐标原点, 力的大小与质点到坐标原点的距离成比例. 设此点依逆时针方向描绘椭圆 $\frac{x^{2}}{a^{2}}+\frac{y^{2}}{b^{2}}=1$ 在第一象限那一段, 求弹性力所做的功.

\item 设在半径为 $R$ 的球上分布着某种物质, 其密度函数 $\rho(\boldsymbol{p})=a\|\boldsymbol{p}\|$, 其中 $a>0$ 为常数. 计算球的质量.

\item 设有半径为 $a>0$ 的球, 球心是 $(0,0,a)$, 密度函数为 $\rho(\boldsymbol{p})=\frac{k}{p}$, 其中 $k$ 为常数, $p=\|\boldsymbol{p}\|$. 计算球的质心.

\item 设有半径为 $R$ 的球, 密度 $\rho$ 为常数 (不妨设 $\rho=1$). 求其引力场. 并说明:

  \begin{enumerate}[label=(\arabic*),leftmargin=2.6em,itemsep=0.15\baselineskip,topsep=0.2\baselineskip]

  \item 对任何一点所产生的引力指向球心;

  \item 对球外一点所产生的引力, 等于在球心上放置一个质量为 $\frac{4}{3}\pi R^{3}$ 的质点对该点所产生的引力, 即犹如球的质量全部集中在球心上;

  \item 对球内一点 $\boldsymbol{p}$ 所产生的引力, 等于半径为 $l$ (点 $\boldsymbol{p}$ 到球心的距离)的球体对点 $\boldsymbol{p}$ 所产生的引力, 犹如球面上一切点若它们到球心的距离大于点 $\boldsymbol{p}$ 到球心的距离时, 它们对 $\boldsymbol{p}$ 点的引力不作贡献.

  \end{enumerate}

\item 设向量场 $\boldsymbol{F}=(x^{2}-2yz,y^{2}-2xz,z^{2}-2xy)$. 证明: $\boldsymbol{F}$ 为有势场, 并求出 $\boldsymbol{F}$ 的势函数.

\item 求有心场 $\boldsymbol{F}$ 的势函数.

\item 一条横截面为半圆形水渠, 求水满时水闸门所承受的压力.
\end{exblock}
